\documentclass[10pt]{article}
\usepackage[margin=21mm]{geometry}
\usepackage{amsmath,amssymb,amsthm,booktabs,hyperref}
\newtheorem{theorem}{Theorem}\newtheorem{lemma}[theorem]{Lemma}\newtheorem{proposition}[theorem]{Proposition}
\newcommand{\ii}{\mathrm i_2}\newcommand{\F}{\mathbb F}\newcommand{\GL}{\operatorname{GL}}\newcommand{\GU}{\operatorname{GU}}\newcommand{\Sp}{\operatorname{Sp}}\newcommand{\PGO}{\operatorname{PGO}}
\newcommand{\PGL}{\operatorname{PGL}}\newcommand{\PGU}{\operatorname{PGU}}\newcommand{\PSL}{\operatorname{PSL}}\newcommand{\PSU}{\operatorname{PSU}}
\setlength{\emergencystretch}{2em}
\title{Torus reversers and explicit bounds for the linear and unitary branches}
\author{Human Compression project}\date{}
\begin{document}\maketitle
\noindent This component concerns $G=\PSL_n(q),\PSU_n(q)$, $n\ge4$.
The involution lower bounds are acquired from the distinctly accepted classical
lower component. The standard descriptions of their full automorphism groups,
finite-field forms and classical matrix orders are retained structural inputs.
We use the explicit upper bound of Lawther--Liebeck--Seitz, Proposition~1.3:
$\ii(\operatorname{Aut}S)<2(q^{N+r}+q^{N+r-1})$ for a simple classical group
of algebraic rank $r$ and $N$ positive roots. The exact maximal-overgroup cover stated in the source Table4 is a named
retained classification premise, with its cited Bereczky/Bamberg--Penttila
provenance and the acquired author erratum. Their full primary texts were
not obtained; their original proofs have not been independently verified.
The arithmetic below does not reprove this classification. This component
recovers only the linear and unitary positive branches, not all of L045.

\paragraph{The chosen tori.}
Write $q=p^f$. In the linear case let $X$ generate the norm-one Singer subgroup
of $\operatorname{SL}_n(q)$ and let $x$ be its projective image. In the unitary
case with odd $n$, use the determinant-one subgroup of the irreducible unitary
Singer torus. For even $n$, use a generator on a nondegenerate $(n-1)$-space,
and its inverse determinant on the complementary line. Put $J=\PGL_n(q)$ or
$\PGU_n(q)$, respectively. The resulting projective torus orders are
\[
 T_L=\frac{q^n-1}{q-1},\qquad
 T_U=\begin{cases}(q^n+1)/(q+1)&n\text{ odd},\\q^{n-1}+1&n\text{ even}.
 \end{cases}
\]
The element orders are these orders divided by $d=\gcd(n,q-1)$ or
$\gcd(n,q+1)$, respectively. Here $\GU$ denotes the unitary isometry group.

\begin{lemma}\label{regular}
For these elements, $C_J(x)=T$. There is no inner-diagonal inverter of $x$.
Every involutory automorphism inverting $x$ lies in the single pure graph
coset of $J$. Consequently their number is at most $|T|$.
\end{lemma}
\begin{proof}
In the linear case $M=|x|$ satisfies
\[
 M\ge\frac{q^n-1}{(q-1)^2}>q^{n/2}+1,
\]
because $q^{n/2}-1\ge q^2-1>(q-1)^2$. A projective centralizer lift
normalizes the field algebra generated by $X$, and so acts on an eigenvalue
by $q^j$, $0\le j<n$. Projective equality implies $M\mid q^j-1$.
For $j\ne0$ this is impossible, since the gcd with $q^n-1$ is at most
$q^{n/2}-1$. Thus the lift is field multiplication and the centralizer is
exactly $T$. An inner inverter would instead give $M\mid q^j+1$.
Squaring gives the same contradiction unless $n$ is even and $j=n/2$;
that case is also excluded by the displayed bound.

The quotient $\operatorname{Aut}(G)/J$ is $C_f\times C_2$, with graph
factor inverse transpose. An involutory reverser having nonzero field
projection has $f$ even and field exponent $f/2$. Eigenvalue comparison
then gives $M\mid p^{f/2+fj}\pm1$. Squaring implies
$M\mid q^{\gcd(n,2j+1)}-1$. A proper gcd is at most $n/2$. Otherwise
$n$ is odd and $2j+1=n$, and the original divisor is at most
$q^{n/2}+1$. Both alternatives contradict the bound. The only remaining
coset is the pure graph coset.

For odd unitary $n$, set $M=(q^n+1)/(d(q+1))$. For $0<j<n$, put
$g=\gcd(n,j)$; then
$\gcd(q^n+1,q^{2j}-1)=q^g+1<M$. Indeed, if $g=1$, $n\ge5$ and
$q^n+1>(q+1)^3$. Otherwise $g\ge3$, $n\ge3g$, and
$q^{2g}-q^g+1>(q+1)^2$. The same eigenvalue argument proves regularity.
For even $n$, the unique one-dimensional eigensummand forces a scalar in
projective commutation to be one, and the irreducible complementary summand
then gives the stated torus centralizer. In both cases the normalizer of
this torus modulo the torus has odd order ($n$ or $n-1$): on its irreducible
summand the normalizer acts by the cyclic field automorphisms. It cannot
induce inversion of $x$. The unitary quotient $\operatorname{Aut}(G)/J$
is cyclic of order $2f$; its sole nonidentity element of order two is the
$q$-Frobenius, or graph, coset. Finally two inverters in the same coset differ
by an element centralizing $x$. This bounds the number by $|T|$; no assertion
that all inverters are conjugate is needed.
\end{proof}

\begin{lemma}\label{fixed}
For every involutory inverter $\alpha$, $C_G(\alpha)$ embeds in
$\PGO_n^\epsilon(q)$ when $q$ is odd, for some form sign, and in the
isometry group $J_n(q)$ of a nonalternating symmetric bilinear form when
$q$ is even. For odd $n$ the latter group is $\Sp_{n-1}(q)$.
\end{lemma}
\begin{proof}
For $\PSL$, the graph form $B$ obeys $X^tB=zBX$ for a scalar $z$.
Similarity between $X$ and $zX$ gives $M\mid q^j-1$, so Lemma~\ref{regular}
forces $z=1$. A bilinear form making field multiplication self-adjoint has
form $B(u,v)=\operatorname{Tr}(cuv)$: write $B(u,v)=\ell(uv)$ and use the
nondegenerate trace pairing. It is symmetric and, in characteristic two,
nonalternating, since $u\mapsto\operatorname{Tr}(cu^2)$ is nonzero.
The projective fixed group of the graph automorphism is its projective
similitude group.

For $\PSU$, take the hermitian form $h$ linear in its first variable, and
represent $\alpha$ by a $q$-semilinear anti-isometry $a$. Its square is $cI$.
Commutation with this square gives $c\in\F_q$, and the isometry condition
gives $c^{q+1}=1$, hence $c=\pm1$. The bilinear form $B(v,w)=h(v,aw)$
satisfies $B(w,v)=cB(v,w)$. The inverter relation makes $X$ self-adjoint
up to a scalar. As above that scalar is one: for odd $n$ use the
$q^{2j}$ eigenvalue comparison in Lemma~\ref{regular}; for even $n$ use
the unique eigenline. Distinct primary summands are $B$-orthogonal because
$X$ is self-adjoint. Its odd-dimensional irreducible summand has a trace
form, which is symmetric and nonalternating in characteristic two.
Thus $c=-1$ is impossible in odd characteristic, and $a^2=1$ in all cases.

Elementary finite-field descent gives $V=\F_{q^2}\otimes W$, where $W$ is
the fixed space of $a$; $h|_W$ is a nondegenerate symmetric $\F_q$ form,
nonalternating in characteristic two. A projective centralizer lift $g$
satisfies $gag^{-1}=ta$ with $t^{q+1}=1$. Choose $u$ with $u^{q-1}=t$;
then $ug$ commutes with $a$ and restricts to a similitude on $W$.
Conversely a similitude on $W$ can be rescaled to a hermitian isometry,
since the norm map is onto. Thus the ambient fixed group is $\PGO(h|_W)$.
In characteristic two every multiplier has a unique square root, so
projective similitudes identify with $J_n(q)$. In odd dimension its
canonical distinguished vector has nonzero length, and its perpendicular
complement is symplectic, proving $J_n=\Sp_{n-1}$.
\end{proof}

\paragraph{The even-dimensional nonalternating group.}
For $n=2m$ in characteristic two choose a basis $z,U,w$ with
$B(z,w)=B(w,w)=1$, $U$ symplectic of dimension $n-2$, and $z$ characterized
by $B(v,v)=B(v,z)^2$. Every isometry has unique parameters
$A\in\Sp(U)$, $v\in U$, $c\in\F_q$:
\[
 z\mapsto z,\quad u\mapsto Au+B(Au,v)z,\quad w\mapsto w+v+cz.
\]
These formulas preserve all pairings and describe every isometry by passing
to $z^\perp/\langle z\rangle$. Squaring gives the identity exactly when
$A^2=1$ and $Av=v$; the remaining coefficient $B(v,Av)$ is then zero.
It follows, including the identity image, that
\begin{equation}\label{Jcount}
 \ii(J_n)=q^{n-1}-1+\sum_{A\in I_2(\Sp_{n-2})}q^{\dim\operatorname{Fix}(A)+1}
 \le q^{n-1}-1+q^{n-2}\ii(\Sp_{n-2}).
\end{equation}
For $n=4$, $\ii(\operatorname{SL}_2(q))=q^2-1$ and every nonidentity
involution has a one-dimensional fixed space, giving
$\ii(J_4)=q^4+q^3-q^2-1$.

\paragraph{Dimension four in odd characteristic.}
The determinant model on $2\times2$ matrices gives
$\PGO_4^+(q)\simeq\PGL_2(q)\wr C_2$: the connected transformations
are $X\mapsto AXB^{-1}$, and adjugation exchanges the factors up to inner
automorphisms. The norm model on hermitian $2\times2$ matrices gives
$\PGO_4^-(q)\simeq\PGL_2(q^2)\rtimes\langle\sigma\rangle$, with
$\sigma$ the field involution. Both maps are faithful modulo the displayed
scalar kernels and have the full projective similitude orders from the
ordinary orthogonal order formula. Since $\ii(\PGL_2(q))=q^2$, the plus
involution count is $(q^2+1)^2-1+|\PGL_2(q)|$. In the minus case every
outer involution is conjugate to $\sigma$: rescale its semilinear lift to
square one and apply finite-field descent. Its class has size
$|\PGL_2(q^2)|/|\PGL_2(q)|=q(q^2+1)$. Thus the two counts are
$q^4+q^3+2q^2-q$ and $q^4+q^3+q$, respectively.

\paragraph{An explicit uniform bound.}
Combining these formulas with LLS Proposition~1.3 gives
$\ii(C_G(\alpha))\le\mu_n(q)$, where
\[
\begin{array}{c|l}
n,q&\mu_n(q)\\\hline
4,\ q\text{ odd}&q^4+q^3+2q^2-q\\
4,\ q\text{ even}&q^4+q^3-q^2-1\\
n\ge5\text{ odd}&2(q^{(n^2-1)/4}+q^{(n^2-1)/4-1})\\
n=2m\ge6,\ q\text{ odd}&2(q^{m^2}+q^{m^2-1})\\
n=2m\ge6,\ q\text{ even}&q^{2m-1}-1+2q^{2m-2}(q^{m^2-m}+q^{m^2-m-1}).
\end{array}
\]
For odd $q$ and even $n\ge6$, the faithful embedding of projective
similitudes in $\operatorname{Aut}(\operatorname{P}\Omega_n^\epsilon)$
uses only that a matrix centralizing the natural derived group is scalar;
$D_3=A_3$ is included. The sole nonsimple symplectic group in these bounds
is $\Sp_4(2)\simeq S_6$, whose $75$ involutions satisfy the stated numerical
upper bound. No use of an automorphism theorem for a nonsimple group is made.

\paragraph{Explicit retained maximal-overgroup cover.}
Retain the published source Table4 assertion that every proper subgroup
containing the chosen $x$ lies
in one of the following actual proper subgroups: in $\PSL_n$, one
extension-field group $(\GL_{n/r}(q^r).r)\cap G$ for each prime $r\mid n$;
in odd-dimensional $\PSU_n$, the corresponding unitary groups for each
prime $r\mid n$; in even-dimensional $\PSU_n$, the single nondegenerate
$(n-1)+1$ stabilizer. The notation means the natural projective image,
with its actual determinant condition. Uniqueness of the extension-field
structures containing the Singer algebra explains ``one'' for each $r$.
Unlike the source's definition of $\mathcal M(x)$, this covering family
\emph{includes} the torus normalizer when $r=n$. For composite degrees the
normalizer preserves the unique prime-degree subfields and is already
contained in an extension-field member. The even-unitary normalizer preserves
the unique eigenline and its perpendicular complement. The element and projective orders here agree exactly with source Table3
in these two families, so the retained cover applies to these actual tori.
Completeness of the cover, with the source's excluded small parameters,
is the named borrowed classification; it is not inferred from these
containment facts or a numerical sample. The acquired Bamberg erratum
adds the odd-prime extension-field groups in odd unitary dimension, already
present in the source's Table4.

\begin{lemma}\label{coverbound}
Under that covering hypothesis, the following explicit number $E_n(q)$
bounds its total involutions. Put $a_k=(k^2+k-2)/2$.
For even-unitary $n$, take
\[
 E_n^U=s\{2(q^{a_{n-1}}+q^{a_{n-1}-1})+1\}-1,
 \qquad s=\begin{cases}1&q\text{ even},\\2&q\text{ odd}.\end{cases}
\]
In the other cases sum, over prime $r\mid n$, the quantities
\[
 E_{n,r}=\begin{cases}
 0&r=n,\\
 (q+1)\{2(q^{2a_{n/2}}+q^{2a_{n/2}-2})+1\}-1&r=2,\\
 s\{2(q^{r a_{n/r}}+q^{r(a_{n/r}-1)})+1\}-1&r\text{ odd},\ r<n.
 \end{cases}
\]
(The $r=2$ case occurs only for the linear family.)
\end{lemma}
\begin{proof}
The map from the linear part of an extension-field subgroup to the smaller
projective group has cyclic scalar kernel. Every fibre of a central cyclic
kernel contains at most two involutions, or at most one when the kernel has
odd order; include the identity image and subtract its identity lift.
For odd $r$, all involutions lie in this linear part, and LLS applies to
the smaller simple group over $q^r$. For $r=2$ the image lies in
$\PGL_{n/2}(q^2).2\le\operatorname{Aut}(\PSL_{n/2}(q^2))$; its kernel
has order at most $q+1$, so use that entire fibre bound instead of assuming
centrality in the field coset. For even-unitary $n$, projection of the
$(n-1)+1$ stabilizer to $\PGU_{n-1}(q)$ has cyclic central kernel, giving
the first formula. If $r=n$ is prime, $n$ is odd, and the projective torus
and its order-$n$ normalizer quotient both have odd order; hence the term
is zero. Every use here is in a simple smaller degree, apart from the
already excluded $\PSU_3(2)$ case. These are upper bounds for actual
subgroups, not assumed equalities between abstract group names.
\end{proof}

\begin{proposition}\label{arithmetic}
For all parameters outside the source's two finite exception lists,
\[
 \ii(G)>|T|\mu_n(q)+E_n(q).
\]
Thus, relative to the stated retained covering classification, those linear and
unitary parameters admit the generating pair required by L045.
\end{proposition}
\begin{proof}
Here are explicit estimates paying the unbounded rank range. For odd
$n\ge7$, $|T|\mu_n\le6q^{(n^2-1)/4+n-1}$. At $n=7$ the covering sum
is zero; the ratio to $q^{(n^2-1)/2}$ is at most $6/64<1/4$.
For odd $n\ge9$, each prime term is at most
$6q^{n^2/6+n/2}$, and there are at most $n$ terms. The sum of the two
ratios is therefore bounded by
\[
6\,2^{-(n-1)(n-3)/4}+6n\,2^{-(2n^2-3n-3)/6}<1/4.
\]
At $n=9$ this is $6/4096+54/2^{22}<1/4$; both summands decrease as
$n$ increases by two (the second exponent increases by at least twelve).
For $n=5$, the sum is zero and, when $q\ge4$,
$|T|\mu_5<(10/3)q^{10}<q^{12}/4$.

For $n=2m\ge8$, the reversing term is at most
$8q^{m^2+3m-3}$. The linear covering sum is at most
$6nq^{m^2+m-1}$, giving ratios at most
\[
8\,2^{-(m^2-3m+3)}+12m\,2^{-(m^2-m+1)}<1/4;
\]
at $m=4$ their sum is $1/16+48/8192$, and both decrease.
For the unitary covering sum in even characteristic the ratio is at most
$3\,2^{-(m+1)}$; adding the reversing bound gives at most
$1/16+3/32<1/4$ at $m=4$, again decreasing.
In odd characteristic the sharper reversing ratio is
$4q^{-(m-1)^2}$ and the covering ratio is at most $6q^{-(m+1)}$;
the sum is already less than $1/4$ at $m=4,q=3$ and decreases.

For $n=6,q\ge4$ the reversing term is less than
$4q^{15}+(4/3)q^{10}$, while either covering sum is less than $6q^{14}$.
The normalized sum is at most $4/q^3+6/q^4+4/(3q^8)<1/4$ at $q=4$
and decreases. For $n=4,q\ge7$, $\mu_4\le(58/49)q^4$,
$|T|<(7/6)q^3$ and either covering sum is less than $5q^5$.
The normalized sum is at most $29/(21q)+5/q^3<1/4$ at $q=7$.
For $n=4,q=5$ the exact accepted class lower bounds and these upper
expressions are, respectively,
\[
\begin{array}{c|r|r}
G&\text{class lower bound}&|T|\mu_4+E_4\\\hline
\PSL_4(5)&251875&131825\\
\PSU_4(5)&170625&115171.
\end{array}
\]
The unused small cases $n=4,q=2,3,4$ and $n=5,6,q=2,3$ are all included
in the separately certified finite exceptions. The acquired lower theorem
$\ii(G)>q^{\lfloor n^2/2\rfloor}/4$ now proves the assertion.

Finally, an involution $t$ outside all cover members generates $G$ with $x$.
If an automorphism inverted $x$ and fixed $t$, its square would fix both
generators and hence be the identity. Thus any such obstruction is already
one of the at most $|T|$ involutory inverters, whose fixed involutions have
been counted. Choosing an involution outside their union gives the required
pair.
\end{proof}

\paragraph{Attribution and limits.}
The source is Leemans--Liebeck, arXiv1603.07713v3, Section4 and Tables3--5.
The new direct torus/descent proofs avoid its small-field application of
Seitz Theorem2.8. The exact radical count and scalar-fibre bounds pay
previously implicit interfaces. Root suggested organizing reversers as
$T$-torsors over outer cosets; the present single-coset exclusions and
linear/unitary estimates are this author's deductions. The unavailable full
classification texts are not represented as acquired or reread, and the source's
incorrect unitary degree-four half bound is not used. The finite35 component
and classical lower theorem retain their separate source and reviewer roles.
\begin{thebibliography}{9}
\bibitem{LL} D. Leemans and M. W. Liebeck, \emph{Chiral polyhedra and finite
simple groups}, arXiv:1603.07713v3, Section4, Tables3--5.
\bibitem{LLS} R. Lawther, M. W. Liebeck and G. M. Seitz,
\emph{Fixed point ratios in actions of finite exceptional groups of Lie type},
Pacific J. Math. 205 (2002), Proposition1.3.
\bibitem{Ber} A. Bereczky, \emph{Maximal overgroups of Singer elements in
classical groups}, J. Algebra 234 (2000), 187--206. Classification retained
through the exact statement in [1], not an acquired full-text proof.
\bibitem{BP} J. Bamberg and T. Penttila, \emph{Overgroups of cyclic Sylow
subgroups of linear groups}, Comm. Algebra 36 (2008), 2503--2543.
Classification retained through [1]; author correction read at
\url{https://johnbamberg.github.io/errata.html}.
\end{thebibliography}
\end{document}

